\specialchapter{术语索引}

绝对平坦（环——）：第 169 页，习题 16。在 n 次零调（复形——）：第 26 页。x-进（拓扑——）：第 159 页。增广代数：第 196 页，习题 27。微分分次代数：第 183 页，习题 18。上升（次数——）：第 23 页。相伴（类——）于正合列：第 117 页。代数的增广：第 196 页，习题 27。自内射（环）：第 171 页，习题 26。双复形：第 174 页，习题 11。上有界、下有界：第 24 页。典范（内射分解——）：第 52 页。典范（自由分解——）：第 50 页。欧拉-庞加莱特征：第 40 页。胞腔（复形——）：第 32 页。戈洛德-沙法列维奇（——定理）：第 192 页，习题 16。五（——引理）：第 7 页。万有系数（——公式）：第 78、99 页。余生成（模——）：第 18 页。左凝聚、右凝聚（环——）：第 181 页，习题 11。凝聚、伪凝聚（模——）：第 181 页，习题 10。群的上同调：第 111 页。上同调（维数——）：第 206 页，习题 28。上同调的：第 194 页，习题 23。余诱导（模——）：第 190 页，习题 12。交换（——同构）：第 71 页。完全切割（族——）：第 157 页。复形：第 23 页。科斯居尔复形：第 147、154 页。复合（——积）：第 113 页。复形态射的锥：第 36 页。锥形（单纯形概形——）：第 177 页，习题 19。联络：第 45 页。收敛（谱序列——）：第 175 页，习题 14。余正则（序列——）：第 159 页。复形态射的柱：第 36 页。戴德金（——环）：第 141、203 页，习题 12。德拉姆（——复形）：第 43 页。下降（次数——）：第 23 页。交换图：第 1 页。微分：第 24 页。微分（代数——）：第 183 页，习题 18。

\begin{Shaded}
\begin{Highlighting}[]
\NormalTok{外微分：第 44 页。}
\NormalTok{上同调维数：第 206 页，习题 28。}
\NormalTok{同调维数：第 138 页。}
\NormalTok{内射维数：第 202 页，习题 4。}
\NormalTok{平坦维数：第 202 页，习题 7。}
\NormalTok{投射维数：第 134 页。}
\NormalTok{可除（模——）：第 17 页。}
\NormalTok{对偶化（模——）：第 172 页，习题 30。}

\NormalTok{可除（模——）：第 17 页。}
\NormalTok{欧拉{-}庞加莱（\textquotesingle{}——的特征）：第 40 页。}
\NormalTok{扩张（——模\textquotesingle{}）：第 81、86 页。}
\NormalTok{一个代数的扩张\textquotesingle{}：第 197 页，习题 1。}
\NormalTok{一个李代数的扩张\textquotesingle{}：第 198 页，习题 4。}

\NormalTok{忠实平坦（模——）：第 169 页，习题 9。}
\NormalTok{微分形式：第 44 页。}

\NormalTok{戈尔迪（——的定理）：第 171 页，习题 25。}
\NormalTok{戈洛德{-}沙法列维奇（——的定理）：第 192 页，习题 16。}

\NormalTok{希尔伯特（——的合冲定理）：第 146 页。}
\NormalTok{同调（类\textquotesingle{}——，模\textquotesingle{}——）：第 25 页。}
\NormalTok{同调（\textquotesingle{}的正合列）：第 30 页。}
\NormalTok{群同调：第 111 页。}
\NormalTok{同调（维数——）：第 138 页。}
\NormalTok{同调关系：第 26 页。}
\NormalTok{同调的（循环——）：第 25 页。}
\NormalTok{同态（复形\textquotesingle{}——）：第 81 页。}
\NormalTok{同伦于零（复形——）：第 34 页。}
\NormalTok{同伦（态射——）：第 32 页。}
\NormalTok{同伦的（单纯地——）：第 177 页，习题 19。}
\NormalTok{同伦：第 32 页。}
\NormalTok{同伦等价：第 33 页。}

\NormalTok{诱导（模——）：第 190 页，习题 12。}
\NormalTok{内射（模——）：第 15 页。}
\NormalTok{内射（维数——）：第 202 页，习题 4。}
\NormalTok{迭代（连接同态——）：第 127、131、132 页。}

\NormalTok{科斯居尔（——的复形）：第 147、154 页。}
\NormalTok{Künneth（——公式）：第 79 页。}

\NormalTok{连接（——同态）：第 29 页。}
\NormalTok{连接（——同态）的扩张模：第 90、92 页。}
\NormalTok{连接（——同态）的挠积：第 72 页。}
\NormalTok{复形的\textquotesingle{}长度：第 24 页。}
\NormalTok{长度（\textquotesingle{}一个表示的）：第 180 页，习题 6。}

\NormalTok{极小（内射分解——）：第 182 页，习题 13。}
\NormalTok{极小（投射分解——）：第 54 页。}
\NormalTok{复形态射：第 25 页。}

\NormalTok{右边为零，左边为零：第 24 页。}
\end{Highlighting}
\end{Shaded}

\begin{Shaded}
\begin{Highlighting}[]
\NormalTok{平坦（绝对——）：第 169 页，习题 16。}
\NormalTok{平坦（忠实——）：第 169 页，习题 9。}
\NormalTok{平坦（维数——）：第 202 页，习题 7。}
\NormalTok{平坦（模——）：第 8 页。}
\NormalTok{庞加莱（特征\textquotesingle{}欧拉{-}——）：第 40 页。}
\NormalTok{庞加莱多项式：第 41 页。}
\NormalTok{前复形：第 174 页，习题 13。}
\NormalTok{表示，有限表示：第 10 页。}
\NormalTok{长度为 n 的表示：第 180 页，习题 6。}
\NormalTok{复合积：第 113 页，}
\NormalTok{挠积：第 61、67 页。}
\NormalTok{投射（维数——）：第 134 页。}
\NormalTok{普吕弗（环——）：第 203 页，习题 11。}
\NormalTok{伪{-}凝聚（模——）：第 181 页，习题 10。}

\NormalTok{正则（局部环——）：第 208 页，习题 12。}
\NormalTok{正则（序列——）：第 158 页。}
\NormalTok{正则（谱序列——）：第 175 页，习题 14。}
\NormalTok{相对投射、相对内射：第 170 页，习题 19。}
\NormalTok{分解：第 48 页。}
\NormalTok{典范内射分解：第 52 页。}
\NormalTok{典范自由分解：第 50 页。}
\NormalTok{分次分解：第 56 页。}
\NormalTok{极小投射分解：第 54 页。}
\NormalTok{一个复形的\textquotesingle{}投射、内射分解：第 183 页，习题 17。}

\NormalTok{单纯形概形：第 177 页，习题 19。}
\NormalTok{分裂（完全——）：第 35 页。}
\NormalTok{分裂：第 35 页。}
\NormalTok{切割（完全——）：第 157 页。}
\NormalTok{蛇（——图）：第 3 页。}
\NormalTok{单纯形：第 178 页，习题 19。}
\NormalTok{单纯（——概形）：第 177 页，习题 19。}
\NormalTok{单纯同伦：第 177 页，习题 19。}
\NormalTok{奇异（——同调）：第 31 页。}
\NormalTok{谱（——序列）：第 175 页，习题 14。}
\NormalTok{标准（——分解）：第 57 页。}
\NormalTok{平稳（——谱序列）：第 175 页，习题 14。}
\NormalTok{同调的\textquotesingle{}正合列：第 30 页。}
\NormalTok{谱序列：第 175 页，习题 14。}
\NormalTok{合冲（——定理）：第 146 页。}

\NormalTok{两个复形的张量（——积）：第 61 页。}
\NormalTok{挠（——积）：第 61、67 页。}
\NormalTok{Tor{-}维数：第 202 页，习题 8。}
\NormalTok{x{-}进拓扑：第 159 页。}
\NormalTok{（p{-}次）平移：第 27 页。}
\NormalTok{转移：第 191 页，习题 14。}
\NormalTok{平凡（上同调——）：第 192 页，习题 18。}

\NormalTok{万有（——系数）：第 78、99 页。}
\end{Highlighting}
\end{Shaded}
